\section[Join and fusion of a principal $H$-comodule algebra with $H$]{Join and fusion of a principal $\boldsymbol{H}$-comodule algebra with $\boldsymbol{H}$}

Since a diagonal coaction is not in general an algebra homomorphism,
to obtain an equivariant version of our noncommutative fusion and join constructions,
we need to modify  Def\/inition~\ref{gjoin}
in the spirit of~\eqref{5},~\eqref{6}.
To avoid analytical complications that
are tackled in~\cite{ddhw}, we also need to change the setting from $C^*$-algebraic to algebraic. In particular,
we replace $\otimes_{\min}$ by~$\otimes$.

\begin{Definition}\label{pjoin}
Let $\delta\colon P\to P\otimes H$ be a right coaction making $P$ a comodule algebra, and
let~$C$ be a unital $C^*$-algebra equipped with two $C^*$-ideals $J_1$ and~$ J_2$
such that the direct sum
\begin{gather*}
\pi_1\oplus\pi_2\colon \ C\longrightarrow (C/J_1) \oplus (C/J_2)
\end{gather*}
of the canonical surjections $\pi_i\colon C\to C/J_i$, $i\in \{1,2\}$, is surjective. We call the algebra
\begin{gather*}
P \underset{\pi}{\overset{\delta}{\circledast}} H :=
\big\{x\in C \otimes P \otimes H \,
\big|\, \\
\hphantom{P \underset{\pi}{\overset{\delta}{\circledast}} H :=\big\{}{}
(\pi_1\otimes \operatorname{id})(x) \in (C/J_1)\otimes \delta (P) ,\,
(\pi_2\otimes \operatorname{id})(x) \in (C/J_2) \otimes \mathbb{C} \otimes H\big\}
\end{gather*}
the \emph{equivariant fusion} algebra of $P$ and~$H$ over~$C$.
When $C:=C([0,1])$  and
$\pi_1:=\mathrm{ev}_1$, $\pi_2:=\mathrm{ev}_0$,
we call $P \circledast^\delta_\pi H$ the \emph{equivariant join} algebra of $P$ and~$H$,
and denote it by $P \circledast^\delta H$.
\end{Definition}

\begin{Lemma}
The coaction $\operatorname{id}\otimes\operatorname{id}\otimes\Delta\colon
C\otimes P\otimes H\to C\otimes P\otimes H\otimes H$
restricts and corestricts to $\Delta_{P\circledast{}^\delta_\pi H}\colon P{\circledast}^\delta_\pi   H
\to (P{\circledast}^\delta_\pi H)\otimes H$ making $P{\circledast}^\delta_\pi H$
a right $H$-comodule algebra.
\end{Lemma}

\begin{proof}
Note f\/irst that for $i\in\{1,2\}$, we have
\begin{gather*}
\big((\pi_i\otimes\operatorname{id}\otimes\operatorname{id})\circ(\operatorname{id}\otimes\operatorname{id}\otimes\Delta)\big)\!\bigg(\sum_j f_j\otimes p_j\otimes h_j\bigg)
=(\operatorname{id}\otimes\operatorname{id}\otimes\Delta)\!\bigg(\sum_j \pi_i(f_j) \otimes p_j\otimes h_j\bigg).
\end{gather*}
Assume now that $\sum_j f_j\otimes p_j\otimes h_j\in P\circledast^\delta_\pi H$.
 Then the above tensor belongs to
\begin{gather*}
(C/J_1)\otimes \big((\operatorname{id}\otimes\Delta)(\delta(P))\big)=
(C/J_1)\otimes \big((\delta\otimes\operatorname{id})(\delta(P))\big)
\subseteq(C/J_1)\otimes \delta(P)\otimes H
\end{gather*}
for $i=1$ and to $(C/J_2)\otimes \mathbb{C}\otimes H\otimes H$ for~$i=2$.
\end{proof}

The main result of this paper is:
\begin{Theorem}\label{main}
Let $C$ be a unital $C^*$-algebra equipped with two $C^*$-ideals $J_1$ and $ J_2$
such that the direct sum
$
\pi_1\oplus\pi_2\colon C\rightarrow (C/J_1) \oplus (C/J_2)
$
of the canonical surjections $\pi_i\colon C\to C/J_i$, $i\in \{1,2\}$, is surjective.
Then, for any principal right $H$-comodule algebra~$P$,
 the equivariant fusion right $H$-comodule algebra $P{\circledast}^\delta_\pi H$ is principal.
\end{Theorem}

\begin{proof}
Note f\/irst that, due to  the surjectivity of $\pi_1\oplus\pi_2$, we can choose $x\in C$  such that
\begin{gather*}
(\pi_1\oplus\pi_2)(x)=(1,0).
\end{gather*}
Since $\{0, 1\}\subseteq \operatorname{spec}(x^*x)$, we can
 take a continuous function $f\colon \operatorname{spec}(x^*x)\to [0,1]$
such that $f(0)=0$ and $f(1)=1$, and def\/ine $t:=f(x^*x)\in C$.
Then there exist elements $\sqrt{t},\sqrt{1-t}\in C$ enjoying the properties
$\pi_2(\sqrt{t})=0=\pi_1(\sqrt{1-t})$.

Next, let
$\ell \colon H{\to} P\otimes P$  be a strong connection on $P$.
Then the bicolinearity of~$\ell$ implies that the linear map
\begin{gather*}
\tilde\ell \colon \  H \longrightarrow \big( C \otimes P \otimes H\big) \otimes \big( C \otimes P \otimes H\big),\\
\tilde\ell(h):= {\sqrt t} \otimes  \ell (h_{(2)})^{\langle 1\rangle}\otimes S(h_{(1)})
\otimes {\sqrt t} \otimes  \ell (h_{(2)})^{\langle 2\rangle}  \otimes h_{(3)}
\\
\hphantom{\tilde\ell(h):=}{}+ \sqrt{1-t} \otimes 1 \otimes S(h_{(1)}) \otimes \sqrt{1-t}
\otimes 1\otimes h_{(2)},
\end{gather*}
corestricts to $(P\circledast^\delta_\pi H)\otimes  (P\circledast^\delta_\pi H)$.
Indeed,
\begin{gather*}
 (\pi_1\otimes\operatorname{id}\otimes\operatorname{id})\otimes(\operatorname{id}\otimes\operatorname{id}\otimes\operatorname{id})(\widetilde{\ell}(h))\\
 \qquad{} =
\big([1]\otimes\delta\big(\ell(h_{(1)})^{\langle 1\rangle}\big)\big)\otimes
\big(\sqrt{t}\otimes \ell(h_{(1)})^{\langle 2\rangle}\otimes h_{(2)}\big)
\in \big(C/J_1\otimes\delta(P)\big)\otimes\big(C\otimes P\otimes H\big),\\
 (\pi_2\otimes\operatorname{id}\otimes\operatorname{id})\otimes(\operatorname{id}\otimes\operatorname{id}\otimes\operatorname{id})(\widetilde{\ell}(h))\\
 \qquad{} =
\big([1]\otimes 1\otimes S(h_{(1)}\big)\otimes
\big(\sqrt{1-t}\otimes 1\otimes h_{(2)}\big)
\in \big(C/J_2\otimes\mathbb{C}\otimes H\big)\otimes\big(C\otimes P\otimes H\big).
\end{gather*}
Much in the same way,
\begin{gather*}
 (\operatorname{id}\otimes\operatorname{id}\otimes\operatorname{id})\otimes(\pi_1\otimes\operatorname{id}\otimes\operatorname{id})(\widetilde{\ell}(h))
\in \big(C\otimes P\otimes H\big)\otimes\big(C/J_1\otimes\delta(P)\big),\\
 (\operatorname{id}\otimes\operatorname{id}\otimes\operatorname{id})\otimes(\pi_2\otimes\operatorname{id}\otimes\operatorname{id})(\widetilde{\ell}(h))
\in\big(C\otimes P\otimes H\big) \otimes\big(C/J_2\otimes\mathbb{C}\otimes H\big).
\end{gather*}

The bicolinearity of the thus corestricted $\tilde\ell$ is evident. Finally, taking
advantage of the right colinearity and the splitting property of~$\ell$, we check that
\begin{gather*}
 \widetilde{\operatorname{can}_{P\circledast^\delta_\pi H}}(\tilde\ell (h)) =
 t \otimes  \ell (h_{(2)})^{\langle 1\rangle}\ell (h_{(2)})^{\langle 2\rangle}
\otimes S(h_{(1)})  h_{(3)}\otimes h_{(4)}\\
\hphantom{\widetilde{\operatorname{can}_{P\circledast^\delta_\pi H}}(\tilde\ell (h)) =}{}
+ (1-t) \otimes 1 \otimes S(h_{(1)})h_{(2)} \otimes h_{(3)}
 \\
\hphantom{\widetilde{\operatorname{can}_{P\circledast^\delta_\pi H}}(\tilde\ell (h))}{} =  t \otimes  \ell (h_{(2)})^{\langle 1\rangle}\ell (h_{(2)})^{\langle 2\rangle}\!{}_{(0)} \otimes S(h_{(1)})
\ell (h_{(2)})^{\langle 2\rangle}\!{}_{(1)}\otimes h_{(3)}
+ (1-t) \otimes 1 \otimes 1 \otimes h
 \\
\hphantom{\widetilde{\operatorname{can}_{P\circledast^\delta_\pi H}}(\tilde\ell (h))}{} =  t \otimes  1 \otimes S(h_{(1)})  h_{(2)}\otimes h_{(3)}
+ (1-t) \otimes 1 \otimes 1 \otimes h
 \\
\hphantom{\widetilde{\operatorname{can}_{P\circledast^\delta_\pi H}}(\tilde\ell (h))}{}
 =  1 \otimes  1 \otimes 1 \otimes h.
\end{gather*}
Now the claim follows from Lemma~\ref{general}.
\end{proof}
